\documentclass[10pt, a4paper]{article}
\usepackage{preamble}

\title{Data Science and Statistical Modelling II}
\author{Luke Phillips}
\date{January 2026}

\begin{document}

\maketitle

\newpage

\tableofcontents

\newpage

\section{Introduction}

$X_i \sim \mathrm{UK}(\mu, \sigma ^ 2)$,
$\mathrm{UK}$ is an unknown distribution,
this goes to $\bar{X} \sim N\left(\mu, \frac{\sigma ^ 2}{n}\right)$ as $n \to \infty$ by CLT.
Unbiased estimate of $\sigma$ is $s = \sqrt{\frac{1}{n - 1}\sum_{i = 1}^{n}(x_i - \bar{x}) ^ 2}$.
The,
$\frac{s}{\sqrt{n}}$ is the standard error of the sample mean.
$n$ small,
$s$ is a poor estimation,
$\bar{X}$ is not Normal.

If assume Normality,
then for unknown $\sigma$,
\[
\frac{\bar{X} - \mu}{\frac{s}{\sqrt{n}}} \sim t_{n - 1}.
\]
What about other parameters?
I.e. median?
Variance?

Data $x_1, \dotsc, x_n$.
Hypothesis test:
$H_0 : \mu = \mu_0$,
$H_1 : \mu \neq \mu_0$.
Choose test statistic such that it is
"extreme" if null is wrong,
not "extreme" if null is ok.
$t = h(x_1, \dotsc, x_n)$.
Need the sampling distribution of test statistic when $H_0$ is true.
$T = h(X_1, \dotsc, X_n)$.
Ask:
how unusual
(is how extreme)
is $t$ under the distribution of $T$ when $H_0$ is true?

Measure:
with $p$-value,,
one-sided:
$\P(T \geq t \mid H_0\text{ true})$ or $\P(T \leq t\mid H_0\text{ true})$.
$2$-sided:
$\P(T \leq -|t| \cup T \geq |t| \mid H_0\text{ true})$.

\section{Monte Carlo Testing}
Idea:
Use simulation instead of calculating the distribution of $T$ when $H_0$ is true.
Data $\{x_1, \dotsc, x_n\}$,
$X_i \sim F(\cdot\mid \theta)$.
$H_0: \theta = \theta_0$,
$H_1 : \theta > \theta_0$,
$t_{\text{obs}} = t = h(x_1, \dotsc, x_n)$.
Monte Carlo testing:
Repeat for $j \in \{1, \dotsc, N\}$
\begin{enumerate}[label = (\roman*)]
    \item Simulate $n$ observations $\{z_1, \dotsc, z_n\}$ for $F(\cdot \mid \theta_0)$.

    \item Calculate $t_j = h(z_1, \dotsc, z_n)$.
\end{enumerate}
At the end we have $\{t_1, \dotsc, t_N\}$.
Then,
$p$-value is
(approximately)
the empirical average $\P(T > t_{\text{obs}} \mid H_0\text{ true}) \approx \frac{1}{N}\sum\mathbb{1}\{t_j > t_\text{obs}\}$,
where
\[
\mathbb{1}\{A\} = \begin{cases}
    1 &\text{if $A$ true} \\
    0 &\text{if $A$ false}.
\end{cases}
\]

Intuition:
Is what we actually saw $t_{\text{obs}}$ unusual compared to many
"parallel universes"
(universes -
$(t_1, \dotsc, t_N)$)
in which $H_0$ really is true?

What problems remain?
\begin{enumerate}[label = (\arabic*)]
    \item Would rather have a standard error and confidence interval.

    \item Still had to assume something for $F(\cdot\mid \theta)$.
\end{enumerate}

\textbf{Bootstrap}:
Data $\vec{x} = (x_1, \dotsc, x_n)$.
Statistic of interest,
$S(\cdot)$.

To construct a Bootstrap estimate of the standard error,
we simulate $B$ samples of size $n$ with replacement from $\{x_1, \dotsc, x_n\}$ giving:
\begin{align*}
    \{x_1 ^ {*1}, \dotsc, x_n ^ {*1}\} &= \vec{x} ^ {*1} \to S(\vec{x} ^ {*1}) \\
    &\vdots \\
    \{x_1 ^ {*B}, \dotsc, x_n ^ {*B}\} &= \vec{x} ^ {*B} \to S(\vec{x} ^ {*B}).
\end{align*}
We take $S(\vec{x})$ as the estimator and compute the variance of the estimate as:
\[
\widehat{\Var}(S(\cdot)) = \frac{1}{B - 1}\sum_{b = 1}^{B}(S(\vec{x} ^ {*b}) - \bar{S} ^ {*}) ^ 2.
\]

Why does bootstrap work?

\begin{definition}
    For a random variable $X$ taking values in $\R$,
    the cumulative distribution function
    (cdf),
    $F : \R \to [0, 1]$,
    is defined by:
    \[
    F(x) \coloneqq \P(X \leq x),\quad \forall x \in \R.
    \]
\end{definition}

\begin{definition}
    Given a sample of IID data,
    $(x_1, \dotsc, x_n)$,
    we define the empirical cdf
    (ecdf)
    to be:
    \[
    \hat{F}(x) \coloneqq \frac{1}{n}\sum_{i = 1}^{n}\mathbb{1}\{x_i \leq x\},\qquad \forall x \in \R.
    \]
\end{definition}

\textit{Note:
this is a valid cdf}
\[
\lim_{x \to -\infty}\hat{F}(x) = 0,\qquad \lim_{x \to \infty}\hat{F}(x) = 1
\]
\[
\forall x' < x, \hat{F}(x') \leq \hat{F}(x)
\]
\[
\forall x \in \R, \hat{F}(x) = \hat{F}(x ^ {+})
\]
$\hat{F}(x ^ {+}) = \lim_{t \to x ^ {+}}\hat{F}(t)$.


\begin{theorem}[Glivenko-Cantelli]
    Let $X_1, \dotsc, X_n$ be a random sample taken from some probability distribution with cdf $F(\cdot)$.
    Then,
    \[
    \sup_{x \in \R}|\hat{F}(x) - F(x)| \to 0\text{ as } n \to \infty.
    \]
\end{theorem}

\begin{lemma}
    Given data $\vec{x} = (x_1, \dotsc, x_n)$,
    then sampling uniformly at random from $\vec{x}$ is equivalent to sampling from the probability distribution defined by the ecdf,
    $\hat{F}(\cdot)$,
    constructed using $\vec{x}$.
\end{lemma}

\begin{enumerate}[label = (\alph*)]
    \item The ecdf $\vec{x}$ approximates the true cdf.

    \item Bootstrap is equivalent to $n$ simulation,
    for the ecdf.

    \item Bootstrap is like fitting ecdf,
    then simulation how the statistic $S(\cdot)$ behaves for samples of size $n$.
\end{enumerate}

The ecdf defines $n$ discrete random variables with pmf:
\[
p(y) = \P(Y = y) = \begin{cases}
    \frac{1}{n} &\text{if } y \in \{x_1, \dotsc, x_n\} \\
    0 &\text{otherwise},
\end{cases}
\]
Calculate moments exactly
\begin{align*}
    \E(Y) &= \sum_{y \in \mathcal{X}}y\cdot p(y) \\
    &= \sum_{i = 1}^{n}x_i\P(Y = x_i) \\
    &= \sum_{i = 1}^{n}x_i\cdot\frac{1}{n} \\
    &= \frac{1}{n}\sum_{i = 1}^{n}x_i \\
    &= \bar{x} \\
    \Var(Y) &= \sum_{y \in \mathcal{X}}(y - \E(Y)) ^ 2p(y) \\
    &= \sum_{i = 1}^{n}(x_i - \bar{x}) ^ 2\cdot\frac{1}{n} \\
    &= \frac{1}{n}\sum_{i = 1}^{n}(x_i - \bar{x}) ^ 2 \\
    &= \frac{n - 1}{n}\cdot\frac{1}{n - 1}\sum_{i = 1}^{n}(x_i - \bar{x}) ^ 2 \\
    &= \frac{n - 1}{n}S_{X} ^ 2.
\end{align*}
Similarly,
we can show $\{Y_1, \dotsc, Y_m\}$ iid draws from the ecdf,
\[
\E(\bar{Y}) = \bar{x},\qquad \Var(\bar{Y}) = \frac{n - 1}{n}\cdot \frac{S_{X} ^ 2}{m}
\]
where $\bar{Y} = \frac{1}{m}\sum_{j = 1}^{m}Y_j$ for $Y_j \sim \hat{F}(\cdot)$.
Hence,
$\sqrt{\Var(\bar{Y})}$ is usual estimate of standard error of mean with correction factor.

Useful but only for $\bar{Y}$,
not so much for the median,
variance,
but we still simulate.

Bootstrap has issues when there is missing data
(e.g. when not all variables are collected for all observations).
Time-series data
(e.g. stock prices,
when strong serial dependence and assumption of independence is not valid).
Censored data
(e.g. data on failure times of engineered system,
could run experiment and observe failures but other units work at the end of the experiment).

Finite population case.
E.g. surveys of finite population of individuals.
Let $N$ be the population size.
$\mu, \sigma ^ 2$ are the population parameters,
mean and variance.
$n$ is the survey size.
$\bar{x}, s ^ 2$ are sample statistics,
mean and variance.
$f \coloneqq \frac{n}{N}$ is the sampling fraction.

If $f > 0.1$
(rough rule of thumb),
then cannot ignore the finite population structure of the problem.
Now,
$\bar{X} = \frac{1}{n}\sum_{i = 1}^{n}X_i$ is still an unbiased estimator for $\mu$.
But,
$\Var(\bar{x})$ is not $\frac{\sigma ^ 2}{n}$ any more.

\begin{theorem}
    Given a finite population size $N$,
    the sample mean based on a sample of size $n$ drawn uniformly without replacement has:
    \[
    \Var(\bar{X}) = \left(\frac{N - n}{N - 1}\right)\frac{\sigma ^ 2}{n}.
    \]

    \begin{proof}
        \begin{align*}
            \Var(\bar{X}) &= \Var\left(\frac{1}{n}\sum_{i = 1}^{n}X_i\right) \\
            &= \frac{1}{n ^ 2}\Var\left(\sum_{i = 1}^{n}X_i\right) \\
            &= \frac{1}{n ^ 2}\left(\sum_{i = 1}^{n}\Var(X_i) + \sum_{i \neq j}\underbrace{\Cov(X_i, X_j)}_{c}\right) \\
            &= \frac{1}{n ^ 2}\left(n\sigma ^ 2 + \sum_{i \neq j}c\right) \\
            &= \frac{1}{n ^ 2}\left(n\sigma ^ 2 + n(n - 1)c\right) \\
            &= \frac{1}{n}(\sigma ^ 2 + (n - 1)c).
        \end{align*}
        If we sample the whole population,
        then $\bar{x} = \mu$,
        $\Var(\bar{X}) = 0$ so
        \[
        0 = \frac{1}{N}(\sigma ^ 2 + (N - 1)c) \implies c = -\frac{\sigma ^ 2}{N - 1}.
        \]
        Hence,
        \[
        \Var(\bar{X}) = \frac{1}{n}\left(\sigma ^2  + (n - 1)\left(-\frac{\sigma ^ 2}{N - 1}\right)\right) = \frac{\sigma ^ 2}{n}\left(\frac{N - n}{N - 1}\right).
        \]
    \end{proof}
\end{theorem}


\end{document}